% Point to the main file
\documentclass[../Main.tex]{subfiles}

% Start the document
\begin{document}

\begin{alphabet}
    
    \item
    {
        Since the paintings can all be in the same room, each painting essentially has 12 options to choose from each, so this turns out to be

        \begin{align}
            \underbrace{12}_{P1} \cdot \underbrace{12}_{P2} \cdot \ldots \cdot \underbrace{12}_{P9} = \boxed{12^9}
        \end{align}
    }

    \item 
    {
        In words, the museum curator would essentially be doing:

        \begin{enumerate}
            \item choose 9 rooms out of 12 for each painting
            \item permute 9 paintings
        \end{enumerate}

        Mathematically, this is

        \begin{align*}
            {12 \choose 9} \cdot \prescript{}{9}{P}_9 &= \boxed{{12 \choose 9} \cdot 9!}
        \end{align*}
    }

    \item
    {
        We can think of this as the following (exactly 1 painting is in the smallest room)

        \begin{center}
            1 in smallest = 0 in biggest and 1 in smallest + at least 1 in biggest and 1 in smallest
        \end{center}

        So

        \begin{center}
            at least 1 in biggest and 1 in smallest = 1 in smallest - 0 in biggest and 1 in smallest
        \end{center}

        For 0 in smallest, each painting just as 11 options (12 - 1 = 11) to exclude the smallest

        So we must choose 1 painting out of 9 to place in the smallest, then the rest has 11 choices for rooms

        \begin{align*}
            \text{1 in smallest} = {9 \choose 1} \cdot \underbrace{1}_{P1} \cdot \underbrace{11}_{P2} \cdot \underbrace{11}_{P3} \cdot \underbrace{11}_{P4} \cdot \underbrace{11}_{P5} \cdot \underbrace{11}_{P6} \cdot \underbrace{11}_{P7} \cdot \underbrace{11}_{P8} \cdot \underbrace{11}_{P8} &= 9 \cdot 11^8
        \end{align*}

        For 0 in biggest (with 1 in smallest), one painting must be in the smallest room, and NONE can be in the biggest.
        This means we must choose one painting out of 9 to be in the smallest, and the rest can be in any of the 10 rooms (can't be in biggest nor smallest, so $12 - 2 = 10$)

        \begin{align*}
            \text{0 in biggest and 1 in smallest} = {9 \choose 1} \cdot \underbrace{1}_{P1} \cdot \underbrace{10}_{P2} \cdot \underbrace{10}_{P3} \cdot \underbrace{10}_{P4} \cdot \underbrace{10}_{P5} \cdot \underbrace{10}_{P6} \cdot \underbrace{10}_{P7} \cdot \underbrace{10}_{P8} \cdot \underbrace{10}_{P9} &= 9 \cdot 10^8
        \end{align*}

        So now we can just plug it into our equation

        \begin{align*}
            \text{at least 1 in biggest and 1 in smallest} &= 9 \cdot 11^8 - 9 \cdot 10^8 \\
            &= \boxed{9(11^8 - 10^8)}
        \end{align*}
    }

    \item
    {
        So our contraint for this one is that A, B, and C must all be in the same room. Because they're all attached to each other, we can just think
        of all three as one unit, so the total number of paintings we have is essentially 7.

        The question then just becomes how many ways can we arrange 7 paintings in 12 rooms?

        Just like the other parts, each painting has 12 options for what room to go in

        \begin{align*}
            \underbrace{12}_{ABC} \cdot \underbrace{12}_{P2} \cdot \underbrace{12}_{P3} \cdot \underbrace{12}_{P4} \cdot \underbrace{12}_{P5} \cdot \underbrace{12}_{P6} \cdot \underbrace{12}_{P7} = \boxed{12^7}
        \end{align*}
    }

    \item
    {
        Somewhat similar to the last one, we can group the 9 paintings into 3 groups.

        Let's define how many ways we can form groups of 3

        \begin{enumerate}
            \item choose 3 out of 9 to be in one group
            \item choose 3 out of 6 to be in another group
            \item choose 3 out of 3 to be in the last group
        \end{enumerate}

        So the number arrangements for groups is given by

        \begin{align*}
            \text{\# of group arrangements} &= {9 \choose 3} \cdot {6 \choose 3} \cdot {3 \choose 3} \\
            &= {9 \choose 3} \cdot {6 \choose 3}
        \end{align*}

        so now we have our groups, we can answer the second part of the question

        \begin{enumerate}
            \item choose 3 out of 12 rooms
        \end{enumerate}

        Which is just

        \begin{align*}
            {12 \choose 3}
        \end{align*}

        So the total arrangements is just

        \begin{center}
            \# of ways to arrange 3 even groups from 9 $\cdot$ \# of ways to choose 3 rooms out of 12
        \end{center}

        \begin{align*}
            \boxed{ {9 \choose 3} \cdot {6 \choose 3} \cdot {12 \choose 3} }
        \end{align*}
    }

\end{alphabet}

% End the document
\end{document}